\documentclass[11pt]{article}

\usepackage[T1]{fontenc}
\usepackage{lmodern}

\usepackage[margin=0.85in]{geometry}
\usepackage{amsmath,amssymb}
\usepackage{booktabs}
\usepackage{microtype}
\usepackage{xcolor}
\usepackage[hidelinks]{hyperref}

\definecolor{navy}{HTML}{142a5e}
\definecolor{accent}{HTML}{c07a1c}
\definecolor{lightbg}{HTML}{eef2f8}
\definecolor{softnavy}{HTML}{2a4a8f}

\setlength{\fboxsep}{9pt}
\setlength{\parindent}{0pt}
\setlength{\parskip}{0.55em}

\hypersetup{
  colorlinks=true,
  urlcolor=accent,
  linkcolor=navy,
  pdfauthor={Rob Mitchell},
  pdftitle={PACLAND --- a new proof about the landscape of everything}
}

% ---- section heading: navy title + amber rule ----
\newcommand{\sechead}[1]{%
  \vspace{1.3em}%
  \noindent{\Large\bfseries\color{navy}#1}\par
  \vspace{0.22em}%
  \noindent\textcolor{accent}{\rule{\linewidth}{1.2pt}}\par
  \vspace{0.65em}%
}

% ---- light callout box with an accent title ----
\newcommand{\callout}[2]{%
  \vspace{0.45em}\noindent
  \fcolorbox{navy}{lightbg}{%
    \parbox{\dimexpr\linewidth-2\fboxsep-2\fboxrule\relax}{%
      {\small\bfseries\color{navy}#1}\par\vspace{0.28em}{\small #2}%
    }%
  }\vspace{0.45em}%
}

% ---- one stat cell in the number strip ----
\newcommand{\statbox}[2]{%
  \colorbox{lightbg}{%
    \parbox{\dimexpr0.30\linewidth-2\fboxsep\relax}{%
      \centering{\color{navy}\fontsize{24}{28}\selectfont\bfseries #1}\\[0.22em]
      {\color{navy}\footnotesize #2}%
    }%
  }%
}

\begin{document}

% ================= HERO =================
\noindent
\colorbox{navy}{%
  \parbox{\dimexpr\linewidth-2\fboxsep\relax}{%
    \vspace{0.35em}
    \begin{center}
      {\color{white}\fontsize{34}{40}\selectfont\bfseries PACLAND}\\[0.12em]
      {\color{accent}\large The landscape, fully characterized}\\[0.55em]
      {\color{white}\small \emph{``Don't ask a complex system for its values --- ask for its
      constraints, its count, and its deficit. Everything else is covariance, not knowledge.''}
      \par\medskip
      String theory has roughly $10^{500}$ possible universes and no experiment that tells most
      of them apart. The answer is not ``which one is ours'' --- it is a proof that the only
      exact structure is a single constraint, that the free parameters are counted (19), and
      that everything else transforms by a fixed law. A theory of everything predicts
      \emph{relations}, never \emph{values}.}
    \end{center}
    \vspace{0.35em}
  }%
}

% ================= NUMBER STRIP =================
\vspace{0.8em}
\noindent\statbox{$10^{500}$}{possible universes in the landscape}
\hfill
\statbox{$1$}{single exact rule that survives}
\hfill
\statbox{$5$}{independent pillars of mathematics}

% ================= 1. THE QUESTION =================
\sechead{The question nobody could settle}

String theory does not predict one universe. It predicts a \emph{landscape} of about
$10^{500}$ of them --- a one followed by five hundred zeros. Each is a different possible
universe, with its own particles, its own forces, its own numbers. Most look nearly
identical to one another, and no experiment can tell most of them apart.

For two decades physicists hunted for a \emph{fingerprint}: a hidden pattern that points to
our universe and lets us rule the other $10^{500}$ out. The hope was that some shortcut --- a
formula, a map, a clue --- was waiting to be found.

% ================= 2. WHAT THE PAPER PROVED =================
\sechead{What the paper proved}

The shortcut is not there --- at least, not the exact kind. Using a 1939 result by the
mathematician Hermann Weyl, the paper proves that the landscape contains exactly one piece
of exact structure: a \textbf{single rule} that every possible universe must obey (physicists
call it a \emph{constraint}). Everything else --- the strengths of the forces, the masses of
the particles, the cosmological constant --- is, in a precise mathematical sense, \emph{noise}.
It carries no fingerprint of the universe that produced it.

\callout{The one rule (don't worry about the symbols)}{The paper writes the rule as
$N = F^{\top} J H$ --- a simple pairing of two numbers, called the D3 tadpole. The point is
not the notation; the point is that this is the \emph{only} exact thing the landscape
contains.}

The headline, now a theorem: \textbf{the landscape's only exact structure is what it
forbids.} No exact, symmetry-respecting formula can single out our universe --- because no
such formula exists to find.

\callout{What this does --- and does not --- rule out}{It rules out \emph{exact algebraic}
selectors: the clean, deterministic formula that picks our universe out of the pile. It does
\emph{not} rule out clever statistical or approximate methods. The search is far harder than
anyone hoped --- not literally impossible.}

% ================= 3. WHY IT MATTERS =================
\sechead{Why it matters}

\textbf{1. It reframes the problem.} The field spent twenty years trying to find the needle
in a $10^{500}$-haystack. This theorem says the right question was never \emph{which needle?}
but \emph{how do we verify a candidate once we have one?} --- and that is a question we can
actually work on.

\textbf{2. It is falsifiable.} Most ideas in this corner of physics cannot be tested. This
one stakes everything on a single, checkable claim: find one exact, symmetry-respecting
quantity that is not just a repackaging of the one rule, and the whole result falls. That is
how science is supposed to work.

\textbf{3. It stands on five pillars.} Five separate branches of twentieth-century
mathematics --- geometry, algebra, symplectic geometry, topology, and counting --- each
independently force the same conclusion. The result does not rest on one argument; it rests
on five.

% ================= 4. WHERE WE GO FROM HERE =================
\sechead{Where we go from here}

The proof covers the gravitational side --- the ``flux'' that shapes a universe. The next
sector is the one we actually live in: the \textbf{Standard Model} of particle physics. There
the same picture is already emerging --- the electric charges are forced by consistency, the
structure is the same two-type pattern --- and one big, falsifiable question is now open: are
the \textbf{19 free numbers} of the Standard Model (the masses and the couplings) predictable
by any deeper theory, or are they, like the landscape, fundamentally unpredictable?
\emph{Either answer is a real discovery.}

In parallel, the proofs are being re-checked by computer --- using \textbf{Lean}, the system
mathematicians use to verify proofs --- so nobody has to take the claim on faith. And the
whole toolkit is open source, free for anyone to run, check, and build on.

% ================= CTA =================
\vspace{1.0em}
\noindent
\colorbox{navy}{%
  \parbox{\dimexpr\linewidth-2\fboxsep\relax}{%
    \vspace{0.3em}
    \begin{center}
      {\color{accent}\large\bfseries See it \,\textbullet\, Check it \,\textbullet\, Build on it}\\[0.45em]
      {\color{white}\small The full story, live:}\\[0.15em]
      {\color{white}\large\bfseries\url{https://pacland.netlify.app}}\\[0.5em]
      {\color{white}\small The code, the proofs, and the data --- everything reproducible:}\\[0.15em]
      {\color{white}\large\bfseries\url{https://github.com/robmitchellengineering-hub/pacland}}\\[0.5em]
      {\color{white}\small The paper itself:}\\[0.15em]
      {\color{white}\normalsize\emph{Mitchell's Completeness Theorem: no exact invariant
      selector in the flux landscape}}
    \end{center}
    \vspace{0.3em}
  }%
}

\vspace{0.7em}
\begin{center}
  {\footnotesize\color{navy}\textbf{PACLAND} \,\textbullet\, Probably Approximately Correct
  Landscape \,\textbullet\, \emph{a verification discipline, not a physics solver}}
\end{center}

\end{document}
